\section{Introduction}
\label{sec:intro}

\IEEEPARstart{H}{yperspectral} imaging (HSI) sensors acquire continuous and finely divided spectral bands across the electromagnetic spectrum, capturing rich spatial-spectral cubes that characterize the physical properties of the Earth's surface. Consequently, HSI classification has emerged as a cornerstone technology in remote sensing, enabling precise land-cover mapping, precision agriculture, environmental monitoring, mineral and resource exploration, and natural disaster assessment~\cite{he2017deep}. Over the past decade, the rapid advancement of deep learning architectures, particularly convolutional neural networks (CNNs) and spatial-spectral Transformers, has dramatically improved classification performance by extracting hierarchical features from these high-dimensional cubes. However, HSI classification remains highly challenging in practical scenarios. The extremely high dimensionality of HSI data, coupled with variations in solar illumination, atmospheric scattering, sensor calibration, and geographical topography, leads to significant variations in spectral signatures. This creates a severe spectral-spatial distribution discrepancy between images captured in different regions or at different times, known as domain shift, which severely degrades the generalization capability of conventional deep learning models when evaluated on out-of-distribution scenes.

To address the performance drop caused by domain shift, transfer learning techniques have been widely researched. Unsupervised domain adaptation (DA) methods aim to align the feature distributions of the source and target domains by minimizing statistical discrepancies or using adversarial training~\cite{ehsnet2023}. However, DA methods are fundamentally limited by the requirement of accessing target domain samples during the model training phase. In real-world remote sensing applications, target-scene data are often completely unavailable during training, making DA impractical for immediate deployment. To circumvent this constraint, domain generalization (DG) has emerged as a more realistic and robust alternative. Unlike DA, DG aims to learn domain-invariant representations using only source domain data, such that the trained model generalizes directly to unseen target domains without any retraining or target data access~\cite{wang2022domain}. Nevertheless, cross-scene HSI domain generalization is exceptionally difficult. The unique, continuous spectral signatures of remote sensing ground objects are highly sensitive to minor environmental variations, leading to complex, non-linear distribution shifts that are difficult to model using only visual annotations from a single source domain.

To overcome the limitations of single-source HSI domain generalization, several approaches have been developed. These can be broadly categorized into generator-based style randomization methods and generator-free vision-language alignment models. Generator-based approaches, such as progressive domain expansion and style randomization frameworks~\cite{li2022progressive}, focus on simulating diverse target-domain distributions by progressively generating intermediate spectral bands and applying style perturbations to source samples. While effective at broadening the visual distribution coverage, these methods rely solely on visual transformations and are prone to generating unrealistic out-of-distribution samples. On the other hand, recent generator-free visual-language models leverage the cross-modal alignment capabilities of pre-trained models like CLIP~\cite{radford2021learning} to align HSI image patches with class-level textual descriptions in a shared semantic space. While these models capture high-level cognitive context, directly adapting them to high-dimensional remote sensing images remains challenging. Standard visual-language models are designed for three-channel natural images and struggle to extract the continuous spatial-spectral signatures and intricate physical properties of multi-band HSI cubes.

These limitations highlight a significant research gap in current HSI domain generalization. First, conventional progressive generation methods rely exclusively on visual-only style variations, which lack explicit semantic guidance. Without semantic constraints, progressive band synthesis often suffers from semantic inconsistency during domain generation, where the structural identity and class-specific layout of the generated patches are distorted. Second, existing methods exhibit insufficient semantic alignment across domains, as they align visual representations without grounding them in domain-invariant language concepts. Finally, the progressive synthesis of HSI bands suffers from weak semantic conditioning, relying on predefined or purely statistics-based band selection schemes during progressive spectral synthesis rather than dynamically integrating multimodal context to preserve class identity. These issues motivate the need for a framework that grounds the progressive generation and classification process in text-driven vision-language semantics, ensuring that generated target-domain variants preserve structural layouts and semantic consistency.

To address these limitations, we propose a Semantic-Guided Progressive Domain Generalization framework for cross-scene HSI classification. The proposed framework integrates high-level semantic descriptions into the progressive spatial-spectral generation process, establishing a unified vision-language learning paradigm. By utilizing a progressive spectral cascade, the model generates domain-generalized target-style variations of HSI patches while explicitly grounding intermediate feature representations in class-specific textual embeddings. This cross-modal guidance enforces semantic consistency during spectral band synthesis, preventing structural distortions and style-induced semantic drift. Finally, we optimize the network using a joint cross-domain semantic alignment strategy, guiding the classifier to capture transferable representations that generalize directly to unseen target domains.

The proposed framework offers several advantages for robust HSI domain generalization. By grounding progressive generation in class semantic representations, the model prevents semantic drift and preserves class-specific structural layouts during cross-domain feature translation. The semantic fusion mechanism enables dynamic alignment between spatial HSI features and textual description embeddings. Additionally, the adaptive semantic conditioning module dynamically modulates the intermediate features of the generator backbone, guiding the network to synthesize realistic target-domain variations. Finally, the robust optimization strategy ensures that the classification boundaries converge to flat loss regions, which stabilizes the training process and enhances generalization on unseen target distributions.

The primary contributions of this work are summarized as follows:
\begin{enumerate}
    \item We propose a novel cross-modal domain generalization paradigm that incorporates text-driven semantic guidance into a progressive generation framework, establishing a robust link between spatial-spectral features and high-level class concepts to mitigate semantic drift under geographical distribution shifts.
    \item We design a semantic-aware progressive generation mechanism that incorporates cross-modal feature fusion and adaptive conditioning, dynamically modulating intermediate spatial-spectral representations to preserve structural layout consistency across domain shifts.
    \item We develop a robust joint optimization framework based on cross-domain semantic alignment and flat-minima searching, ensuring that the model converges to generalize effectively across unseen target distributions.
    \item Extensive experiments on three public HSI benchmark datasets validate the domain generalization capability of our approach, demonstrating the effectiveness and competitiveness of the proposed framework under diverse cross-scene settings.
\end{enumerate}
